\section{\texorpdfstring{\(L_1\) 正则为何产生稀疏}{L1 正则为何产生稀疏}}
\label{sec:09-Convex-Optimization/07-Applications/02}

上线约束是这样下来的：特征服务每次请求最多取二十个字段，多一个就要多一次远程调用，延迟预算吃不消。手上的候选特征有三千多个。Ridge 跑完，三千多个系数没有一个是零，最小的那个是 \(4\times10^{-7}\)，最大的是 \(0.83\)，中间连成一条光滑的谱，找不到任何一处可以下刀的缝。

按绝对值排序取前二十个？这是把选择重新交回给人：阈值定在哪里没有依据，而且被砍掉的两千九百多个系数所携带的拟合，在砍完之后没有任何机制补回到留下的那二十个身上。更合理的做法是让筛选发生在优化过程内部——让目标函数自己把该丢的系数推到恰好为零，剩下的系数在零的约束下重新拟合。

做法只是换一个惩罚项：

\[
\min_{w}\ \frac12\norm{Xw-y}_2^2+\lambda\norm{w}_1 .
\]

问题仍然是凸的，仍然是最小二乘骨架，唯一的改动是平方换成绝对值。而解的性质变了：\(\lambda\) 调到合适的位置，绝大多数坐标会精确地等于零，不是等于 \(10^{-7}\)。这一节要说明这个零是从哪来的。答案在几何里，不在统计假设里。

\subsection{尖角长在坐标轴上}

先把惩罚形式换成等价的约束形式，因为可行域看得见：

\[
\begin{aligned}
\text{minimize}\quad & \tfrac12\norm{Xw-y}_2^2\\
\text{subject to}\quad & \norm{w}_1\le t .
\end{aligned}
\]

平方损失的等值线是一族同心椭圆，中心在无约束最小二乘解 \(\hat w\)，越往外目标值越大。约束把可行域限制成一个球。最优解在哪里？从中心开始把椭圆一圈圈放大，第一次碰到可行域的那个点就是解——再小的椭圆够不着可行域，再大的椭圆已经越过它。

于是全部问题变成：这一族椭圆第一次碰到球，碰在什么位置。而这取决于球长什么样。

\begin{center}
\begin{tikzpicture}[font=\small,>=stealth,scale=0.92]
  % ---------- 左：L1 球 ----------
  \draw[black!25,->] (-2.4,0) -- (2.95,0) node[right,font=\footnotesize,black!60] {\(w_1\)};
  \draw[black!25,->] (0,-1.6) -- (0,3.3) node[above,font=\footnotesize,black!60] {\(w_2\)};
  \fill[black!8] (1.1,0) -- (0,1.1) -- (-1.1,0) -- (0,-1.1) -- cycle;
  \draw[black!55] (1.1,0) -- (0,1.1) -- (-1.1,0) -- (0,-1.1) -- cycle;
  \begin{scope}[shift={(1.578,1.984)},rotate=20]
    \draw[black!30] (0,0) ellipse (1.14 and 0.51);
    \draw[black!80,thick] (0,0) ellipse (1.9 and 0.85);
    \draw[black!30] (0,0) ellipse (2.47 and 1.105);
  \end{scope}
  \fill[black!70] (1.578,1.984) circle (1.5pt);
  \node[font=\footnotesize,black!65] at (2.3,2.3) {\(\hat w\)};
  \fill[black!85] (0,1.1) circle (2.2pt);
  \node[font=\footnotesize,black!75] at (-1.3,-0.72) {接触点是顶点};
  \node[font=\footnotesize,black!75] at (-1.3,-1.18) {顶点上 \(w_1=0\)};
  \node[font=\footnotesize,black!70] at (1.2,-1.18) {\(\norm{w}_1\le t\)};
  % ---------- 右：L2 球 ----------
  \begin{scope}[shift={(7.2,0)}]
    \draw[black!25,->] (-2.4,0) -- (2.95,0) node[right,font=\footnotesize,black!60] {\(w_1\)};
    \draw[black!25,->] (0,-1.6) -- (0,3.3) node[above,font=\footnotesize,black!60] {\(w_2\)};
    \fill[black!8] (0,0) circle (1.1);
    \draw[black!55] (0,0) circle (1.1);
    \begin{scope}[shift={(1.578,1.984)},rotate=20]
      \draw[black!30] (0,0) ellipse (1.14 and 0.51);
      \draw[black!80,thick] (0,0) ellipse (1.80 and 0.805);
      \draw[black!30] (0,0) ellipse (2.47 and 1.105);
    \end{scope}
    \fill[black!70] (1.578,1.984) circle (1.5pt);
    \node[font=\footnotesize,black!65] at (2.3,2.3) {\(\hat w\)};
    \fill[black!85] (0.293,1.060) circle (2.2pt);
    \node[font=\footnotesize,black!75] at (-1.3,-0.72) {接触点在光滑弧上};
    \node[font=\footnotesize,black!75] at (-1.3,-1.18) {两个分量都非零};
    \node[font=\footnotesize,black!70] at (1.2,-1.18) {\(\norm{w}_2\le t\)};
  \end{scope}
\end{tikzpicture}

\emph{同一族椭圆等值线，同一个中心 \(\hat w\)：碰到菱形时落在顶点上，那里 \(w_1\) 恰好为零；碰到圆时落在一段光滑弧上，两个分量都不为零。}
\end{center}

图上两边的椭圆完全一样，差别只在灰色区域的形状。\(\ell_1\) 球是把正方形旋转 \(45^\circ\) 得到的菱形，四个顶点恰好坐在两条坐标轴上；\(\ell_2\) 球是圆，处处光滑。椭圆胀大时先碰到菱形凸出来的那个尖角，而尖角所在的位置 \(w_1=0\)——一个坐标被精确地置零了。右边的圆没有尖角可碰，接触点落在一段弧上，弧上除了四个孤立的点以外，两个分量都不为零，而椭圆恰好停在那四个点上是零概率事件。

这个对照要读到的东西是：稀疏来自可行域的形状，与真实模型是否稀疏无关。就算生成数据的机制里三千个特征全部有微弱的作用，\(L_1\) 照样会把大部分系数压成零——尖角不问真相，它只是在那里。稀疏是优化结构的产物，是你选择的惩罚项的几何后果，不是从数据里发现的事实。这句话后面还要用到，因为它直接决定了稀疏结果能承担多少解释责任。

推广到 \(d\) 维：\(\ell_1\) 球是有 \(2d\) 个顶点的多胞形，顶点位于坐标轴上；它的低维面也各自落在若干个坐标超平面 \(w_j=0\) 的交上。椭圆碰在 \(k\) 维面上，就意味着 \(d-k\) 个坐标同时为零。维数越高，这些带尖的面在边界上占的比重越大，稀疏就越容易发生。

\subsection{零点的次梯度是一段区间，不是一个数}

几何解释了倾向，还需要一个能逐坐标验算的条件。\(\norm{w}_1\) 在任何一个坐标为零的地方不可微——这正是尖角在函数层面的样子。绝对值函数在零点的次微分是整段区间：

\[
\partial\abs{u}\big|_{u=0}=[-1,1],
\]

含义见\hyperref[sec:09-Convex-Optimization/06-Optimization-Algorithms/03]{``不可微时次梯度为何仍能前进''}。把整个目标的最优性条件写成 \(0\in X^{\trans}(Xw-y)+\lambda\,\partial\norm{w}_1\)，记残差 \(r=y-Xw\)、\(x_j\) 为第 \(j\) 列特征，逐坐标展开就是三条：

\[
\begin{aligned}
w_j>0 &\ \Longrightarrow\ x_j^{\trans}r=\lambda,\\
w_j<0 &\ \Longrightarrow\ x_j^{\trans}r=-\lambda,\\
w_j=0 &\ \Longrightarrow\ \abs{x_j^{\trans}r}\le\lambda .
\end{aligned}
\]

问题是凸的，这三条合起来既必要又充分，于是可以反过来读。前两条说，所有被选中的特征与残差的相关量的绝对值全都等于 \(\lambda\)，一个不多一个不少。第三条是关键：只要某个特征与残差的相关量落在宽度为 \(2\lambda\) 的带子里，取 \(w_j=0\) 就已经满足最优性——次微分那段区间里总找得出一个值把等式配平，不需要把这个坐标从零推开。

换成 \(L_2\) 惩罚，同一个位置的导数是单点 \(\set{0}\)，没有区间可用。此时 \(w_j=0\) 要成为最优，必须 \(x_j^{\trans}r\) 精确等于零。只要相关量还剩下 \(10^{-7}\)，最优解就必须让 \(w_j\) 离开零去把它抵消掉。开头那三千个非零小数，逐个都是这条等式的产物。尖角与光滑的几何差别，在一阶条件里就是``一段区间''与``一个点''的差别。

\subsection{正交设计下最优解就是软阈值}

取一个能算到底的特例。设各列已标准化且互相正交，\(X^{\trans}X=I\)，记 \(z=X^{\trans}y\)，也就是逐特征做一元回归得到的系数。此时

\[
\tfrac12\norm{Xw-y}_2^2=\tfrac12\norm{w-z}_2^2+\text{常数},
\]

目标按坐标完全分离，第 \(j\) 个坐标只需最小化 \(\tfrac12(w_j-z_j)^2+\lambda\abs{w_j}\)。分段求导后

\[
w_j^\star=\sign(z_j)\max\set{\abs{z_j}-\lambda,\ 0}.
\]

这是软阈值算子。它把 \(\lambda\) 的含义讲得很直白：\(\lambda\) 是入选门槛。相关量的绝对值超过门槛的特征活下来，但系数被向零拉近 \(\lambda\)——这一段收缩就是稀疏的价钱，被选中的特征也被压小了，估计因此有偏。没超过门槛的特征直接归零。相比之下 \(L_2\) 的对应算子是 \(w_j^\star=z_j/(1+\lambda)\)，按比例缩小，永远够不到零。

顺手能读出正则化路径的起点。若 \(\lambda\ge\norm{X^{\trans}y}_\infty=\max_j\abs{z_j}\)，所有坐标都满足上一节的零条件，全零解最优。所以 \(\lambda_{\max}\) 是从上往下扫描的自然起点：\(\lambda\) 从 \(\lambda_{\max}\) 开始下降，特征按相关量的大小依次被激活进入模型。实践中扫一条对数等距的 \(\lambda\) 网格并用上一个解热启动，比每个 \(\lambda\) 独立求解快一个数量级。

一般的设计矩阵没有正交性，坐标之间通过 \(X^{\trans}X\) 的非对角元耦合，闭式解消失。但软阈值不会消失，它换了个位置出现——作为近端算子。近端梯度法每一步先沿光滑部分的梯度走一步，再对整个向量施加软阈值；坐标下降每次固定其余坐标只优化一个，子问题恰好退化成上面那个一维问题，解还是软阈值，只是 \(z_j\) 换成扣除其他坐标贡献后的部分残差相关量。两种算法的构造见\hyperref[sec:09-Convex-Optimization/06-Optimization-Algorithms/04]{``近端与坐标方法利用可分结构''}。值得注意的是，正是因为软阈值有闭式解，不可微的 \(L_1\) 反而比很多光滑惩罚更好算——不可微在这里没有变成计算障碍。

\subsection{被选中不等于重要}

现在回到交付场景。Lasso 交出二十个非零系数，报告里能不能写``这二十个是关键因素''？

不能，至少不能只凭这个结果写。设想两列特征几乎完全相同，就像上一篇里那对补录字段。它们对残差的相关量几乎相等，谁进模型都能把这份拟合吃下来，目标函数在``此消彼长''的方向上几乎没有曲率。数据里一点抖动就足以让解从一个尖角滑到另一个尖角。换一次随机划分，进模型的可能是另一列，而两次报告读起来会像在讲两个不同的故事。

更严格地说，当 \(X\) 不满足受限特征值一类的条件时，Lasso 的最优解未必唯一：同一组数据、同一个 \(\lambda\)，可以有多个稀疏模式完全不同的最优解，它们的目标值一模一样。求解器报告的是其中一个，选中哪一个由算法细节和初值决定。这不是实现的缺陷，是问题本身没有把答案定死——和上一篇里平谷上停在哪一点是同一类现象，只不过那里的解集是一条连续的线，这里是一组离散的尖角。

Elastic Net 是常见的折中：

\[
\min_w\ \tfrac12\norm{Xw-y}_2^2+\lambda_1\norm{w}_1+\tfrac{\lambda_2}{2}\norm{w}_2^2 .
\]

\(L_1\) 保住尖角，稀疏机制还在；\(L_2\) 补上 \(\lambda_2\) 的强凸性，解重新唯一，而且高度相关的一组特征倾向于一起进入或一起出局，不再是随机点名。代价是多一个超参数要调，且被选中的系数收缩得更狠。

标准化在这里比在 Ridge 里更要紧。\(L_1\) 对所有坐标用同一个门槛 \(\lambda\)，特征的量纲直接决定它被砍的先后：同一个物理量以毫米计比以米计的系数小三个数量级，会被优先归零。所以先在训练折内标准化，再谈选择；截距不进惩罚项。这两条做错，选出来的就是量纲表而不是特征表。

\subsection{\texorpdfstring{\(L_1\) 是计数函数的凸包络}{L1 是计数函数的凸包络}}

退一步看，最初的诉求其实是``非零系数不超过二十个''，写出来是

\[
\norm{w}_0=\#\set{j:w_j\ne0}\le 20 .
\]

这个函数既不连续也不凸，约束下的最小化是组合搜索，从三千个里挑二十个有 \(\binom{3000}{20}\) 种取法，属于 NP-hard 的一类，见\hyperref[sec:09-Convex-Optimization/09-Complexity-and-Approximation/01]{``P、NP 与 NP-hard：给问题按计算难度分类''}。

\(L_1\) 是它在 \(\norm{w}_\infty\le1\) 上的凸包络——所有落在 \(\norm{w}_0\) 下方的凸函数中最大的那一个。这个身份解释了两件事。它是所有凸下界里最紧的，所以在凸的范围内没有更好的替代品；它毕竟只是下界，对大系数的惩罚随 \(\abs{w_j}\) 线性增长，而 \(\norm{w}_0\) 对已经非零的系数不再追加惩罚，于是 \(L_1\) 会把该留的系数也压小。凸松弛的一般套路见\hyperref[sec:09-Convex-Optimization/09-Complexity-and-Approximation/02]{``凸松弛把 NP-hard 问题变成可解问题''}。

松弛什么时候是紧的，即 \(L_1\) 的解恰好落在真实的支撑集上，取决于设计矩阵的性质。压缩感知给出的受限等距性和不相干性条件是这类回答的代表：列之间足够接近正交，恢复就有保证；列高度相关，保证就失效，退回上一节那种谁进模型都行的局面。

同一套思路可以换掉复杂度的度量而保留结构。若关心的是整组系数同时为零，把组内取 \(\ell_2\) 范数再对组求和，得到 group lasso；若关心的是矩阵的秩，注意秩就是奇异值向量的 \(\norm{\cdot}_0\)，它的凸包络是核范数，即奇异值之和，于是低秩恢复也变成凸问题；若关心相邻系数是否相等，惩罚差分的 \(L_1\)，得到融合 lasso。每一次都是同一个动作：写下真正关心的那个不凸的复杂度，取它的凸包络，再交给近端算法。每一次也都伴随同一个代价：解出来的是替代目标的最优解，它离原目标有多远，取决于问题是否满足那些结构条件。

\subsection{从稀疏的系数到稀疏的样本}

回头清点这一篇的因果链。惩罚项带尖角，尖角落在坐标轴上，等值线先碰到尖角，于是坐标为零；一阶条件里，尖角表现为一段可以吸收相关量的次梯度区间；算法层面，尖角表现为一个有闭式解的软阈值算子。三种语言说的是同一件事，而它们都只谈优化结构，不谈数据的真实分布。

值得追问的是，不可微的尖角是否只能长在惩罚项上。如果把尖角放进损失函数——比如让损失在样本分类正确到一定程度之后就完全归零，之前则线性上升——那么被``归零''的就不是系数，而是样本对解的影响。少数落在转折处的样本决定整条边界，其余样本挪动一下也不改变答案。这正是下一篇\hyperref[sec:09-Convex-Optimization/07-Applications/03]{``最大间隔与支持向量机''}的主题，那里稀疏出现在对偶变量上。
